\subsection{Sprint 3}
Na Sprint 3, as atividades planejadas concentraram-se na finalização da tela do player de áudio que havia sido desenvolvida separadamente e em sua integração ao projeto principal. Essa tarefa era o coração da aplicação, garantindo que os usuários pudessem acessar as meditações semanais com facilidade. A tela foi concluída com sucesso e aceita pela cliente na entrega da sprint, representando um progresso significativo no cumprimento dos objetivos estabelecidos para essa fase do projeto.

\begin{figure}[H]
    \centering
    \small
    \caption{Tela do \textit{player} de áudio integrada ao aplicativo do meditAMAmente}
    \includegraphics[width=0.4\linewidth]{conteudo/3 - ages II/conteudo/figures/tela_do_player_de_audio.png}\\
    \label{fig:medita-player-audio}
    \href{https://tools.ages.pucrs.br/meditamamente/wiki/-/wikis/mockups}{\textit{Fonte: https://tools.ages.pucrs.br/meditamamente/wiki/-/wikis/mockups}}
\end{figure}

Apesar do sucesso na entrega e na finalização da /\textit{task}, a sprint trouxe desafios inesperados, principalmente relacionados à migração da tela do player de áudio para o projeto. Embora a tela parecesse pronta, ela foi desenvolvida em um ambiente isolado e o processo de incorporá-la ao sistema principal revelou dificuldades que não haviam sido previstas. A tela estava mal componentizada e utilizava uma arquitetura significativamente diferente do restante do projeto, o que dificultou sua adaptação. Esses problemas tornaram a tarefa mais demorada e exigiram um esforço adicional para alinhar a tela às diretrizes de arquitetura e melhorar sua organização. Além disso, o processo de integração se mostrou mais confuso do que o esperado, e, devido ao tempo reduzido, acabei precisando solicitar ajuda de outros membros da equipe em cima da hora para garantir a conclusão da tarefa no prazo.

Essa sprint me trouxe lições valiosas para projetos futuros. Um dos aprendizados mais importantes foi a necessidade de não subestimar tarefas, mesmo quando parecem estar adiantadas ou parcialmente concluídas. A avaliação inicial de complexidade pode não considerar detalhes que emergem apenas durante o desenvolvimento ou integração. Desse modo, é necessário evitar a medição de esforços baseada em um julgamento prévio da dificuldade de uma tarefa. Além disso, ficou claro que é essencial iniciar as atividades com antecedência suficiente para que eventuais dificuldades possam ser superadas sem comprometer o prazo, incluindo a possibilidade de buscar apoio quando necessário.

Com a conclusão da sprint 3, o foco do time agora se direciona para o desenvolvimento e a integração das últimas \ac{us}, como a tela de diários da mãe, além da preparação e entrega dos relatórios necessários. Assim, apesar dos desafios enfrentados, essa sprint demonstrou a importância da adaptação, do trabalho em equipe e da busca por soluções colaborativas para superar dificuldades, contribuindo para o crescimento individual e coletivo no projeto.